\section{Appendix B: Rational realization and the upper bound}
\label{sec:rational-realization}

We begin with the representation and primitive arithmetic contracts used by the finite rational construction.

The projection bounds in \cref{obj:lem:projection-control} supply the statistical ingredients for the coordinate error radii. This appendix establishes their finite rational realization: an explicit admissible arithmetic engine, uniform guarantees at the two prescribed precisions, and a measurable interval output. These conclusions connect the exact coordinate expressions in \cref{obj:def:coordinate-estimates} to the literal interval procedure in \cref{obj:def:upper-handle}.

The two precision levels control rank selection and final endpoint enclosure. Each primitive is evaluated by a finite rational recurrence with a public iteration bound on valid represented inputs. The guarantees establish termination and endpoint enclosure under the stated rational computation model, with public bounds on primitive iteration counts.

\subsection{Naming conventions and primitive admissibility}

Under \cref{obj:def:dyadic-name-convention}, a real input is represented by nested closed dyadic intervals containing it, with width at most \(2^{-q}\) at integer precision \(q\ge0\). An admissible Borel naming policy \(\mathsf N\) selects public exponent names before sampling and covariate names as Borel functions of the public exponents and original records. The dyadic-floor policy \(\mathsf N_0\) provides the concrete selection specified there.

The independently fixed engine \(\mathsf E\) satisfies the primitive contracts of \cref{obj:def:arithmetic-engine}. Its transcendental routines enclose the images of rational input intervals, with outward endpoint excess controlled by the dyadic tolerance \(\epsilon_q=2^{-q}\). The cosine routine uses a midpoint enclosure enlarged by half the input width and clipped to \([-1,1]\). Exact rational range operations handle arithmetic, extrema, intersections, and clipping. Domain and sign conditions ensure positive logarithm arguments, nonzero division denominators, positive exponential and power lower endpoints, and nonnegative square-root lower endpoints. Each primitive uses a finite rational recurrence with a prescribed public bound on its iteration count. In the power routine, the bound argument \(b\) denotes a positive integer base.

These requirements specify local arithmetic guarantees. The following statement supplies an explicit engine satisfying them.

\begin{lemmav}{lem:concrete-arithmetic-engine}[Admissibility of the concrete engine]
The explicit rational arithmetic engine \(\mathsf E_0\), consisting of the Machin enclosure for \(\pi\), the cosine interval routine, and the endpoint extensions of the scalar exponential, logarithm, square-root, and integer-base power routines, satisfies all primitive admissibility contracts of \cref{obj:def:arithmetic-engine}.
\end{lemmav}

The witness in \cref{obj:lem:concrete-arithmetic-engine} establishes existence at the primitive level. Its finite rational routines provide one concrete choice for the interval construction. The precision guarantees below apply uniformly to every admissible engine, including this witness.

\subsection{Uniform realization at fixed precision}

Rank selection follows \cref{obj:def:resolutions}: the public exponent intervals are intersected with the prescribed exponent boxes, and the ceilings of the upper target-enclosure endpoints determine the two integer ranks. These ranks are fixed before sampling. Endpoint evaluation then uses precisely those retained integers in the coordinate estimates, variance envelope, error radii, and precision budget of \cref{obj:def:upper-handle}. At that step, the queried exponent and covariate intervals enter the endpoint expression tree in their supplied form.

The next statement establishes admissibility of the concrete naming policy and derives the rank, endpoint, and measurability guarantees from the primitive contracts.

\begin{lemmav}{lem:fixed-budget-realization}[Fixed-budget realization and measurability]
The dyadic-floor naming policy \(\mathsf N_0\) is admissible in the sense of \cref{obj:def:dyadic-name-convention}. Moreover, suppose that:
\begin{itemize}
\item \textbf{(Engine and names.)} \(\mathsf E\) is an admissible rational interval-arithmetic engine in the sense of \cref{obj:def:arithmetic-engine}, and \(\mathsf N\) is an admissible Borel naming policy in the sense of \cref{obj:def:dyadic-name-convention}.
\item \textbf{(Sample size.)} \(n\) is an integer with \(n\ge2\).
\item \textbf{(Public exponents.)} The real smoothness exponents satisfy
\[
0<\beta<\frac14,\qquad \beta<\alpha\le1.
\]
\end{itemize}
For every record tuple \(\mathcal O\in([0,1]\times\{0,1\}\times\{0,1\})^n\), let \(k_C,k_S\) be the ranks selected according to \cref{obj:def:resolutions}, with target resolutions
\[
R_C=n^{2/(2\alpha+2\beta+1)},\qquad
R_S=n^{2/(4\beta+1)}.
\]
The prescribed public budgets are
\[
q_{\mathrm{rank}}(n)=7+3\lceil\log_2 n\rceil,\qquad
q_{\mathrm{end}}(n,k_C,k_S)
=44+\lceil\log_2(n+1)\rceil
 +4\lceil\log_2(K+1)\rceil,\qquad
K=\max\{k_C,k_S\}.
\]
At the rank budget, both rational target enclosures \([a_C,b_C]\) and \([a_S,b_S]\) are successfully returned and satisfy
\[
b_C-a_C\le1,\qquad b_S-a_S\le1,
\]
\[
R_C\le k_C<R_C+2\le3R_C,\qquad
R_S\le k_S<R_S+2\le3R_S.
\]
Here \(b_C,b_S\) denote the upper endpoints of the resolution enclosures.

For the endpoint construction, use \(b_C,b_S\) instead to denote the coordinate error radii
\[
b_C=8k_C^{-(\alpha+\beta)}+\sqrt{20V_n(k_C)},\qquad
b_S=25k_S^{-2\beta}+\sqrt{20V_n(k_S)},
\]
where \(V_n(k)\) is the public variance envelope at rank \(k\). Let \(\widehat C_n,\widehat S_n\) be the numerator and denominator coordinate estimates, and set
\[
c_\pm=\widehat C_n\pm b_C,\qquad
s_\pm=\operatorname{clip}_{[s_0,1]}(\widehat S_n\pm b_S),
\]
where \(s_0\) is the denominator floor. With \(F(c,d)\) denoting the clipped logarithmic inversion of the coordinate ratio, the ideal lower and upper endpoints are
\[
L=\min\{F(c_-,s_-),F(c_-,s_+)\},\qquad
R=\max\{F(c_+,s_-),F(c_+,s_+)\}.
\]
At the endpoint budget, rational enclosures \([L_-,L_+]\) and \([R_-,R_+]\) are successfully returned with
\[
L\in[L_-,L_+],\qquad R\in[R_-,R_+],\qquad
L_+-L_-\le\frac1n,\qquad R_+-R_-\le\frac1n.
\]
The literal output satisfies
\[
[L,R]\subseteq
I^{\mathrm{up}}_{n,\mathsf E,\mathsf N}(\alpha,\beta)(\mathcal O),
\qquad
\operatorname{length}\!\left(
I^{\mathrm{up}}_{n,\mathsf E,\mathsf N}(\alpha,\beta)(\mathcal O)
\right)
\le R-L+\frac2n.
\]
It is a nonempty closed connected interval with rational endpoints in \([-1/2,1/2]\), and its output map, as a function of the record tuple, is Borel measurable.
\end{lemmav}

The rank precision \leanref{sym:q_{\mathrm{rank}}}{\ensuremath{q_{\mathrm{rank}}(n)}} gives an additive rank excess below two for each coordinate. The endpoint precision \leanref{sym:q_{\mathrm{end}}}{\ensuremath{q_{\mathrm{end}}(n,k_C,k_S)}} depends on the sample size and the larger retained rank. The same selected integers therefore govern both the exact statistics and their rational evaluation. Together with the finite recurrence contracts, successful evaluation at these prescribed precisions establishes termination of the construction in \cref{obj:def:upper-handle}.

In the inversion, \(R\) denotes the ideal upper effect endpoint. The rational output uses the outward endpoints \(L_-,R_+\), followed by clipping to the effect envelope. Its samplewise containment of the ideal interval and additional length bound of \(2/n\) hold for every record tuple. The Borel conclusion concerns the resulting map from the original observed records, under each fixed admissible naming policy.

\subsection{Connection to honesty and expected length}

The rational interval \leanref{sym:I^{\mathrm{up}}}{\ensuremath{I^{\mathrm{up}}_{n,\mathsf E,\mathsf N}(\alpha,\beta)}} in \cref{obj:def:upper-handle} combines these realization guarantees with the statistical bounds. The bias and variance bounds in \cref{obj:lem:projection-control} calibrate its coordinate error radii; the denominator floor and odds identity in \cref{obj:lem:observable-odds} supply the joint inversion. \Cref{obj:lem:fixed-budget-realization} establishes measurable outward enclosure of the resulting ideal interval with a controlled length increment.

Under the engine, naming-policy, and exponent conditions of \cref{obj:thm:finite-honest-upper}, the literal procedure has global \(90\%\) coverage over the ambient model for every integer sample size \(n\ge1\). On every effect-radius slice, that result bounds its unconditional expected length by an absolute constant times the diagnostic envelope in \cref{obj:def:diagnostic-envelope}. Numerator uncertainty contributes directly to this envelope, while denominator uncertainty contributes at the scale of the effect radius. At \(n=1\), the specified full effect interval completes the procedure.

Different admissible engines and valid naming policies may select different ranks and return different rational endpoints. For each choice, the realization bounds apply to its own retained ranks and ideal endpoints. The coverage and expected-length guarantees in \cref{obj:thm:finite-honest-upper} share one absolute constant across all such choices. In particular, \cref{obj:lem:concrete-arithmetic-engine,obj:lem:fixed-budget-realization} supplies the explicit engine and dyadic-floor policy for the concrete attaining sequence defined in \cref{obj:def:frontier-handle}.
